\chapter{Laryngectomy}

\section*{Definition and Overview}
Laryngectomy is a major surgical procedure involving the partial or complete removal of the larynx (the voice box). It is primarily indicated for advanced or recurrent laryngeal and hypopharyngeal malignancies, severe laryngeal trauma, or radiation-induced necrosis where the structure and function of the larynx are irreversibly damaged. In a total laryngectomy, the entire larynx, including the thyroid and cricoid cartilages, epiglottis, and hyoid bone, is excised. This completely disconnects the upper digestive tract from the respiratory tract. Consequently, the trachea is redirected and sutured to the skin of the anterior neck, creating a permanent breathing port known as a tracheostoma (or stoma). Because the natural vocal cords are removed, speech is lost, and patients must undergo vocal and respiratory rehabilitation. Partial laryngectomy techniques preserve sections of the larynx, allowing the retention of some natural voice and airway function, though they are only suitable for early-stage or localised disease.

\section*{Epidemiology}
Laryngeal cancer is the principal driver for laryngectomy, representing approximately 1\% of all diagnosed cancers in males globally. In Australia, the age-standardised incidence is approximately 3 to 4 cases per 100,000 people annually. There is a strong gender disparity, with the condition being roughly 8 times more common in men than in women, largely reflecting historic smoking and alcohol consumption patterns. The peak incidence occurs between the ages of 60 and 75. While the overall rate of total laryngectomies has declined over the past few decades due to the adoption of organ-preservation protocols (combining chemotherapy and radiotherapy), it remains an essential salvage procedure for radiation failures and advanced stage IV disease.

\section*{Aetiology and Pathophysiology}
The need for a laryngectomy arises from pathological conditions that destroy the laryngeal cartilage, compromise the airway, or infiltrate adjacent structures.
\begin{itemize}
    \item \textbf{Laryngeal squamous cell carcinoma (SCC):} The primary cause, accounting for over 95\% of all laryngeal malignancies. Chronic exposure to carcinogens leads to dysplastic changes in the squamous epithelium lining the vocal cords (glottis), supraglottis, or subglottis.
    \item \textbf{Tobacco smoking:} The single most significant risk factor. Carcinogens in tobacco smoke induce DNA damage and mutations in tumor suppressor genes (such as TP53), leading to uncontrolled cell proliferation.
    \item \textbf{Alcohol consumption:} Acts synergistically with tobacco smoke. Alcohol functions as a solvent, enhancing the penetration of tobacco-derived carcinogens into the laryngeal mucosa.
    \item \textbf{Chondrosarcoma and other rare tumours:} Non-epithelial malignancies of the laryngeal cartilages (such as the cricoid or thyroid cartilage) that do not respond well to radiotherapy and require surgical resection.
    \text{\item \textbf{Radiation necrosis:}} Chronic, progressive tissue death and chondroradionecrosis of the laryngeal cartilage following high-dose radiotherapy for head and neck cancers, resulting in a non-functional, painful, or obstructed larynx.
    \item \textbf{Severe trauma:} Massive crush injuries or penetrating wounds to the neck that disrupt the skeletal framework of the larynx beyond the possibility of functional reconstruction.
\end{itemize}

\section*{Clinical Features}
The presentation of laryngeal disease necessitating laryngectomy depends heavily on the subsite of the larynx affected (glottis, supraglottis, or subglottis).
\begin{itemize}
    \item \textbf{Persistent hoarseness:} The most common and early sign of glottic tumours, caused by irregular vocal fold vibration. Any hoarseness lasting more than three weeks warrants urgent investigation.
    \item \textbf{Dyspnoea and stridor:} High-pitched, noisy breathing and shortness of breath indicating significant airway narrowing by a tumour or vocal cord fixation.
    \text{\item \textbf{Dysphagia and odynophagia:}} Difficulty and pain when swallowing, often indicating tumor spread into the piriform sinuses, hypopharynx, or base of the tongue.
    \item \textbf{Chronic cough and haemoptysis:} A persistent cough, occasionally productive of blood-stained sputum, arising from ulcerated mucosal tumours.
    \item \textbf{Referred otalgia:} Ear pain, typically unilateral, mediated via the auricular branch of the vagus nerve (Arnold's nerve) due to shared sensory pathways.
    \item \textbf{Cervical lymphadenopathy:} A palpable, firm, non-tender lump in the neck, representing regional lymphatic metastasis.
    \text{\item \textbf{Cachexia and weight loss:}} Significant, unintended weight loss due to a combination of swallowing difficulties and tumour-induced systemic inflammatory state.
\end{itemize}

\section*{Differential Diagnosis}
Before committing to a laryngectomy, clinicians must differentiate malignant laryngeal disease from benign pathologies and other non-laryngeal conditions.
\begin{itemize}
    \item \textbf{Vocal cord nodules and polyps:} Benign, reactive lesions secondary to vocal abuse or chronic irritation, typically resolving with speech therapy or conservative surgical excision.
    \item \textbf{Chronic laryngitis:} Persistent inflammation of the laryngeal mucosa caused by gastro-oesophageal or laryngopharyngeal reflux, chronic smoking, or recurrent infections.
    \item \textbf{Vocal cord paralysis:} Loss of movement in one or both vocal folds due to recurrent laryngeal nerve damage (e.g. from thoracic aortic aneurysm, apical lung tumours, or prior thyroid surgery).
    \item \textbf{Hypopharyngeal and cervical oesophageal cancer:} Tumours arising outside the larynx that present with dysphagia and referred ear pain, but may secondarily invade the larynx.
    \item \textbf{Thyroid carcinoma:} Advanced thyroid malignancies invading the trachea, larynx, or recurrent laryngeal nerves.
    \item \textbf{Laryngeal tuberculosis or sarcoidosis:} Rare granulomatous diseases that can mimic neoplastic lesions on laryngoscopy and require biopsy for differentiation.
\end{itemize}

\section*{When to Seek Medical Care}
An evaluation by an ear, nose, and throat (ENT) specialist is necessary for any patient presenting with persistent hoarseness or throat pain lasting longer than three weeks, unexplained neck lumps, or swallowing difficulties.

\textbf{Call 000 (or local emergency services) for:}
\begin{itemize}
    \item Acute, severe breathing difficulty or stridor (noisy, gasping inhalation).
    \item Inability to swallow saliva or secure the airway.
    \item Severe bleeding from the throat or a post-operative tracheostoma.
    \item Sudden blockage or dislodgement of a laryngectomy tube or stoma in a post-operative patient.
\end{itemize}

\section*{Diagnosis and Investigations}
A comprehensive diagnostic workup is required to confirm the diagnosis, determine the tumor stage, and plan the surgical margins.
\begin{itemize}
    \item \textbf{Flexible laryngoscopy:} An in-clinic endoscopic examination of the larynx and pharynx to assess vocal fold mobility, tumour location, and mucosal appearance.
    \item \textbf{Direct laryngoscopy and biopsy under general anaesthesia:} The definitive test to obtain tissue samples for histopathological confirmation and mapping the exact boundaries of the disease.
    \item \textbf{Contrast-enhanced CT scan of the neck and chest:} Essential to evaluate the cartilage invasion, pre-epiglottic space involvement, cervical lymph node status, and to screen for pulmonary metastases.
    \item \textbf{Magnetic resonance imaging (MRI) of the neck:} Provides superior soft-tissue contrast to delineate submucosal tumor extension, tongue base infiltration, or perineural spread.
    \item \textbf{Positron emission tomography (PET/CT):} Used for staging advanced disease to detect occult regional or distant metastases and to monitor for recurrence.
    \item \textbf{Multidisciplinary pre-operative workup:} Comprehensive assessments by a speech pathologist, dietitian, dentist, and clinical nurse consultant.
\end{itemize}

\section*{Management}
Management is highly complex and requires an integrated multidisciplinary team, including head and neck surgeons, oncologists, speech pathologists, and specialised nurses.
\begin{itemize}
    \item \textbf{Surgical resection (Laryngectomy):} The removal of the larynx. In a total laryngectomy, this is often accompanied by a unilateral or bilateral selective neck dissection to remove regional lymph nodes.
    \item \textbf{Tracheostoma creation:} Direct suturing of the tracheal stump to the skin of the lower neck, establishing the stoma.
    \item \textbf{Voice rehabilitation:} Re-establishing communication. The gold standard is a tracheoesophageal puncture (TEP), where a small tract is created between the trachea and oesophagus to house a one-way silicone voice prosthesis. Other options include an electrolarynx or training in oesophageal speech.
    \item \textbf{Radiotherapy and chemotherapy:} Administered either as adjuvant therapy post-operatively for high-risk features (extracapsular nodal spread, positive margins) or as primary curative therapy.
    \item \textbf{Nutritional support:} Placement of a nasogastric tube (NGT) or percutaneous endoscopic gastrostomy (PEG) tube to maintain nutrition during the early healing phase when oral intake is suspended.
    \item \textbf{Psychosocial support:} Counseling and peer support groups to assist patients with the significant emotional and social adjustment of living with a stoma.
\end{itemize}

\section*{Complications}
Laryngectomy is a major surgical procedure carrying significant risk of early postoperative and long-term complications.
\begin{itemize}
    \item \textbf{Pharyngocutaneous fistula:} A communication between the pharyngeal closure line and the neck skin, allowing saliva to leak into the neck tissues, which can lead to infection and carotid artery erosion.
    \item \textbf{Stoma stenosis:} Narrowing of the tracheostoma due to scar contracture, requiring surgical revision (stomal dilation or stomaplasty) or long-term use of a stoma vent.
    \text{\item \textbf{Mucus plugging and crusting:}} Accumulation of thick, dry tracheal secretions that can block the airway, particularly in the immediate post-operative period.
    \item \textbf{Hypocalcaemia:} Caused by inadvertent damage or removal of the parathyroid glands during laryngectomy and thyroidectomy, requiring calcium and vitamin D supplementation.
    \item \textbf{Hypothyroidism:} Developing months to years post-operatively due to partial or complete thyroidectomy or adjuvant radiation, necessitating lifelong thyroxine replacement.
    \item \textbf{Haemorrhage:} Early post-operative bleeding from major cervical vessels, which is a surgical emergency.
\end{itemize}

\section*{Prognosis and Outcomes}
The prognosis after a laryngectomy depends primarily on the pathological stage of the primary tumour and the presence of nodal metastases. For patients with stage III and IV laryngeal cancer undergoing salvage or primary total laryngectomy, the 5-year overall survival rate ranges from 50\% to 65\%. Patients with localized disease who fail primary radiation and undergo salvage laryngectomy often achieve excellent long-term local control. Voice rehabilitation with a TEP is highly successful, with up to 80\% of patients achieving functional, intelligible speech. However, the physical change of breathing through a neck stoma and the loss of natural voice can have a profound impact on quality of life, requiring ongoing psychological and rehabilitative support.

\section*{Prevention and Self-Care}
Primary prevention revolves around lifestyle modifications, while post-operative self-care is vital for airway maintenance and overall well-being.
\begin{itemize}
    \item \textbf{Smoking and alcohol cessation:} Eliminating tobacco and alcohol significantly reduces the risk of developing secondary primary malignancies in the head and neck region.
    \item \textbf{Tracheostoma humidification:} Since the upper airway no longer filters and humidifies inspired air, patients must use humidifiers, saline instillations, or heat and moisture exchangers (HMEs) to prevent dry trachea and crusting.
    \item \textbf{Stoma cleaning and suctioning:} Daily cleaning of the stoma skin, gentle suctioning of secretions, and regular cleaning of laryngectomy or tracheostomy tubes.
    \item \textbf{Water precautions:} Strict avoidance of water entering the stoma during bathing, showering, or washing. Swimming is strictly contraindicated unless using highly specialised laryngectomy swimming equipment.
    \item \textbf{Emergency identification:} Wearing a medical alert bracelet or carrying a card identifying the patient as a "neck breather" (incapable of mouth-to-mouth resuscitation; emergency oxygen/ventilation must be delivered via the stoma).
\end{itemize}

\section*{Sources and Further Reading}
The following organisations provide further information on laryngectomy and head and neck cancer:
\begin{itemize}
    \item Cancer Council Australia. \textit{Understanding Head and Neck Cancers.} https://www.cancer.org.au/
    \item Laryngectomee Association of New South Wales. \textit{Support and resources for laryngectomees.} http://www.laryngectomee.org.au/
    \item Mayo Clinic. \textit{Laryngeal Cancer: Diagnosis, treatment and recovery.} https://www.mayoclinic.org/
    \item NHS. \textit{Laryngectomy overview and rehabilitation.} https://www.nhs.uk/
    \item Cancer Research UK. \textit{Living with a laryngectomy.} https://www.cancerresearchuk.org/
\end{itemize}
